Given a unitary $U_{\mathrm{diag}}$ expressed in a diagonalizing frame, such as in Eq.~\eqref{eq:diag_unitary}, we show that any Pauli-$Z$ interaction strength $\theta_S$ can be can recovered by linear combinations of phases acquired by computational-basis states under the application of $U_{\mathrm{diag}}$. \\

\noindent
Since $U_{\mathrm{diag}}$ is diagonal in the computational basis, its action on basis states
$\ket{\vec{x}}$, with $\vec{x}\in\{0,1\}^n$, is
\begin{equation}
U_{\mathrm{diag}}\ket{\vec{x}}
=
e^{-i\phi(\vec{x})}\ket{\vec{x}},
\end{equation}
where $\phi(\vec{x})$ denotes the phase acquired by the basis state $\ket{\vec{x}}$. The full unitary can therefore be written as
\begin{equation}
U_{\mathrm{diag}}
=
\sum_{\vec{x}\in\{0,1\}^n}
e^{-i\phi(\vec{x})}
\ket{\vec{x}}\bra{\vec{x}},
\label{eq:diag_unitary_comp_phases}
\end{equation}
so that we write its generator in terms of computational-basis phases, 
\begin{equation}
\Phi
=
\sum_{\vec{x}\in\{0,1\}^n} \phi(\vec{x})\,|\vec{x}\rangle\langle\vec{x}|,
\quad \text{with} \quad
U_{\mathrm{diag}} = e^{i\Phi}.
\label{eq:phi_generator}
\end{equation}
Since $\Phi$ is diagonal in the computational basis, it admits a natural expansion in the diagonal Pauli-$Z$ operator basis
\(Z_S = \prod_{i\in S} Z_i,\,\text{with}\, S \subseteq \{1,\dots,n\}, \)
so that it can also be written as a unique combination of $Z_S$ terms
\begin{equation}
\Phi
=
\sum_{S\subseteq\{1,\dots,n\}}
\phi(S)\, Z_S ,
\label{eq:phase_pauli_expansion}
\end{equation}
where $\phi(S)$ are linear combinations of the coefficients $\phi(\vec{x})$ in Eq.\eqref{eq:phi_generator}. \\
Since Eqs.~\eqref{eq:diag_generator} and \eqref{eq:phase_pauli_expansion} provide two equivalent representations of the same diagonal generator in the Pauli-$Z$ basis, the corresponding expansion coefficients must coincide,
\begin{equation}
\phi(S)=\theta_S.
\label{eq:phi_theta_equivalence}
\end{equation}
The explicit form of the coefficients follows from projection onto the Pauli-$Z$ basis using the trace inner product,
\begin{equation}
\phi(S)
=
2^{-n}\mathrm{Tr}(\Phi Z_S)
=
2^{-n}
\sum_{\vec{x}\in\{0,1\}^n}
\langle\vec{x}|\Phi Z_S|\vec{x}\rangle.
\label{eq:phase_invariant}
\end{equation}
Since computational-basis states are eigenstates of Pauli-$Z$ strings,
\begin{equation}
Z_S|\vec{x}\rangle
=
(-1)^{\sum_{i\in S}x_i}|\vec{x}\rangle,
\end{equation}
the coefficients reduce to
\begin{equation}
\phi(S)
=
2^{-n}
\sum_{\vec{x}\in\{0,1\}^n}
(-1)^{\sum_{i\in S}x_i}
\langle\vec{x}|\Phi|\vec{x}\rangle.
\end{equation}
Using
\(
\langle\vec{x}|\Phi|\vec{x}\rangle=\phi(\vec{x}),
\)
we obtain
\begin{equation}
\phi(S)
=
2^{-n}
\sum_{\vec{x}\in\{0,1\}^n}
(-1)^{\sum_{i\in S}x_i}
\phi(\vec{x}).
\label{eq:phase_coefficients_parity_sum}
\end{equation}

\medskip
\noindent
Thus the coefficients $\theta_S$ can be uniquely obtained as parity-weighted sums of computational-basis phases. Equation~\eqref{eq:phase_coefficients_parity_sum} is the Walsh-Hadamard transform of the computational-basis phase map, equivalently the projection of the diagonal phase generator onto the Pauli-$Z$ basis. We use the coefficients $\phi(S)=\theta_S$ as \emph{interaction coordinates}, since each coefficient directly specifies the interaction contribution supported on the subset $S$.

\medskip
\noindent
\textbf{Remark.}
The Walsh-Hadamard form makes the support resolution explicit: orthogonality of the parity functions separates the Pauli-$Z$ contributions associated with distinct qubit subsets~\cite{welch2014efficient}.\\

\noindent
The interaction coordinates $\phi(S)$~\eqref{eq:phase_coefficients_parity_sum} provide an interaction-resolved coordinate system for unitaries in a diagonalizing frame, as they decompose any diagonalized unitary into $k$-body contributions. The diagonal unitary can therefore be written as
\begin{equation}
U_{\mathrm{diag}}
=
\exp\!\left(-
i\sum_{S\subseteq\{1,\dots,n\}}
\phi(S)\, Z_S
\right),
\label{eq:phi_expansion_u_diag}
\end{equation}
where the coefficients $\phi(S)$ specify the many-body contribution strength supported on the subset $S$.


\subsection{Equivalence Classes}
\label{sec:interaction_equivalence_classes}

\noindent
Since the coefficients $\phi(S)$ provide an interaction-resolved coordinate system through Eq.~\eqref{eq:phi_expansion_u_diag}, selected coordinates can be constrained while the remaining coordinates are left free. This can be expressed conveniently through equivalence relations specifying which interaction coordinates are required to coincide.

For a fixed support subset $S$, two unitaries $U$ and $V$, represented in diagonalized form, are said to be $S$-interaction equivalent if they possess the same coordinate along the interaction direction $Z_S$:
\begin{equation}
\phi_U(S)
=
\phi_V(S)
\quad \mathrm{mod}\; 2\pi.
\end{equation}

More generally, for a chosen collection of supports
\[
\mathcal{A}\subseteq 2^{\{1,\dots,n\}},
\]
we define the interaction-equivalence relation
\begin{equation}
U \sim_{\mathcal{A}} V
\quad\Longleftrightarrow\quad
\phi_U(S)=\phi_V(S)
\; \mathrm{mod}\; 2\pi,
\qquad
\forall S\in\mathcal{A}.
\end{equation}


Thus, each selected set of supports $\mathcal{A}$ determines a corresponding equivalence class on \(U(2^n)\) via the interaction-coordinate space: unitaries belonging to the same equivalence class agree on the specified many-body interaction coordinates while remaining free to differ along coordinates associated with subsets outside $\mathcal{A}$.

In particular, choosing
\[
\mathcal{A}_{\geq2}
=
\left\{
S\subseteq\{1,\dots,n\}
:\ |S|\geq2
\right\}
\]
fixes all nonlocal interaction coordinates while allowing arbitrary one-body coordinates. This recovers an analogue of local equivalence within the interaction-resolved representation: local phases may vary, but the nonlocal interaction-coordinate structure is preserved.\\

\noindent
Conversely, selecting a single support $S$ defines a more flexible equivalence class containing all unitaries sharing the same $S$-body interaction coordinate, independent of all remaining coordinates.\\

\noindent
These equivalence classes are useful for quantum control because they allow objectives to be formulated at different levels of interaction resolution within the same coordinate framework. Rather than matching a full target unitary, one may require only membership in a prescribed interaction-equivalence class, for example by constraining selected interaction coordinates $\phi(S)$ while leaving local or otherwise irrelevant coordinates unconstrained.

% \medskip
% \noindent


